\documentclass[10pt,a4paper]{amsart}
\usepackage[margin=19mm]{geometry}
\usepackage{amsmath,amssymb,mathtools}
\usepackage[scr=rsfs,cal=euler]{mathalfa}
\usepackage{xfrac}
\usepackage[unicode,hidelinks,bookmarks=false]{hyperref}
\newcommand{\ie}{i.e.\ }
\newcommand{\cf}{cf.\ }
\newcommand{\resp}{respectively\ }
\newcommand{\Subsection}{\subsection}
\providecommand{\nobreakdash}{\nobreak\hbox{-}}
\setlength{\emergencystretch}{2em}
\multlinegap0pt
\usepackage{multirow}
\usepackage{amssymb}
\usepackage{xcolor}
\usepackage{booktabs}
\usepackage{algorithm}\usepackage{algpseudocode}
\usepackage{tikz}
\newtheorem{lemma}{Lemma}
\usepackage{cleveref}
\crefname{lemma}{Lemma}{Lemmas}
\newcommand{\closure}[1]{\overline{#1}}
\newcommand{\ctrace}[3]{\gamma^\parallel(#1,#2,#3)}
\newcommand{\dualbasis}{\mathcal{L}}
\newcommand{\polyspace}{\mathcal{V}}
\newcommand{\refel}{R}
\newcommand{\subentity}{E}
\newcommand{\unctrace}[3]{\gamma^\perp(#1,#2,#3)}
\title{DefElement: an encyclopedia of finite element definitions}
\author{Matthew W. Scroggs, Pablo D. Brubeck, Joseph P. Dean, J{\o}rgen S. Dokken, India Marsden}
\date{Source: arXiv:2506.20188v2 (CC BY 4.0)}
\begin{document}
\maketitle

\noindent\emph{Source and licence.} DefElement: an encyclopedia of finite element definitions. Version arXiv:2506.20188v2 (CC BY 4.0), distributed by arXiv under the Creative Commons Attribution 4.0 International licence, http://creativecommons.org/licenses/by/4.0/, as recorded in the arXiv licence metadata for this exact version and shown on the abstract page https://arxiv.org/abs/2506.20188v2. Full licence notice: \texttt{licenses/arxiv-2506.20188-CC-BY-4.0.txt}.\par\smallskip

\noindent\textit{Editorial scope.} Counted unit: the article's lemma lemma:first-c (source0.tex lines 1744--1753). The article's complete original proof of it is delivered with it (source0.tex lines 1754--1803, one \texttt{proof} environment of 50 lines). No mathematics is written, repaired or re-derived: the statement and the proof are the article's own lines, in the article's own notation, and no retained line is substituted or re-flowed.  Retained uncounted as the source-local context the proof needs: lines 1633--1638; lines 1698--1705; lines 1728--1738.  Nothing was omitted from any retained span, so the package carries no omission record.  No retained line contains a citation, so the wrapper carries no bibliography and no bibliography entry.  No retained line contains a figure environment, an image-inclusion command or drawing code, so no image is delivered and the package declares no asset dependency.  Editorial formatting only, in addition: an AMS \texttt{amsart} preamble that restates the article's own declarations and macros used by the retained lines, namely the environments \texttt{lemma} on separate counters, together with the standard packages those lines need.  The retained environments print in this document as Lemma 1 in order of appearance, whereas the article prints the same items with the numbering of its own full document.  The complete native source of the article as distributed in the arXiv e-print for this exact version is bundled unchanged as \texttt{sources/180479/source0.tex}.
\begin{align}
\left.\polyspace\right|_{\subentity}
&=
\unctrace{\polyspace}{\subentity}{\dualbasis}\oplus\ctrace{\polyspace}{\subentity}{\dualbasis}.
\label{oplus}
\end{align}

\begin{lemma}\label{lemma:first-unc}
Let $(\refel,\polyspace,\dualbasis)$ be a reference-mapped finite element.
If $\subentity$ is a sub-entity of $\refel$, then
\[
\unctrace{\polyspace}{\subentity}{\dualbasis} =
\operatorname{span}\left\{\left.\phi_i\right|_{\subentity}\,:\,l_i\not\in\left.\dualbasis\right|_{\closure{\subentity}}\right\}.
\]
\end{lemma}

\begin{lemma}\label{lemma:zeroth}
Let \(\left.\polyspace\right|_{\subentity}:=\operatorname{span}\left\{\left.\phi_0\right|_{\subentity},\dots,\left.\phi_{n-1}\right|_{\subentity}\right\}\).
Given $f\in\left.\polyspace\right|_{\subentity}$, there are unique functions
$f^\parallel\in\ctrace{\polyspace}{\subentity}{\dualbasis}$
and
$f^\perp\in\unctrace{\polyspace}{\subentity}{\dualbasis}$
such that
\[f=f^\perp+f^\parallel.\]
Moreover, the operator $P_{\subentity}:\left.\polyspace\right|_{\subentity}\to\ctrace{\polyspace}{\subentity}{\dualbasis}$
defined by $P_{\subentity}:f\mapsto f^\parallel$ is a linear projection.
\end{lemma}

\begin{lemma}\label{lemma:first-c}
Let $(\refel,\polyspace,\dualbasis)$ be a reference-mapped finite element.
If $\subentity$ is a sub-entity of $\refel$, then
\begin{align*}
\ctrace{\polyspace}{\subentity}{\dualbasis} &=
\operatorname{span}\left\{P_{\subentity}\left(\left.\phi_i\right|_{\subentity}\right)
\,:\,l_i\in\left.\dualbasis\right|_{\closure{\subentity}}\right\},
\end{align*}
where \(P_{\subentity}\) is defined as in \cref{lemma:zeroth}.
\end{lemma}

\begin{proof}
Using \cref{oplus}, we see that
\begin{align*}
\unctrace{\polyspace}{\subentity}{\dualbasis}\oplus\ctrace{\polyspace}{\subentity}{\dualbasis}
&=\left.\polyspace\right|_{\subentity}\\
&=
\operatorname{span}\left\{\left.\phi_0\right|_{\subentity},\dots,\left.\phi_{n-1}\right|_{\subentity}\right\}\\
&=
\operatorname{span}\left(\left\{\left.\phi_i\right|_{\subentity}\,:\,l_i\not\in\left.\dualbasis\right|_{\closure{\subentity}}\right\}
\cup
\left\{\left.\phi_i\right|_{\subentity}\,:\,l_i\in\left.\dualbasis\right|_{\closure{\subentity}}\right\}\right).
\end{align*}
Using \cref{lemma:first-unc}, we therefore see that
\begin{align*}
\unctrace{\polyspace}{\subentity}{\dualbasis}\oplus\ctrace{\polyspace}{\subentity}{\dualbasis}
&=
\unctrace{\polyspace}{\subentity}{\dualbasis}
+
\operatorname{span}\left\{\left.\phi_i\right|_{\subentity}\,:\,l_i\in\left.\dualbasis\right|_{\closure{\subentity}}\right\}.
\end{align*}
By \cref{lemma:zeroth}, we know that for all $l_i\in\left.\dualbasis\right|_{\closure{\subentity}}$,
there exists $g_i\in\unctrace{\polyspace}{\subentity}{\dualbasis}$ such that
\(\left.\phi_i\right|_{\subentity}=g_i+P_{\subentity}\left(\left.\phi_i\right|_{\subentity}
\right)\). It follows that
\begin{equation}
\unctrace{\polyspace}{\subentity}{\dualbasis}\oplus\ctrace{\polyspace}{\subentity}{\dualbasis}
=
\unctrace{\polyspace}{\subentity}{\dualbasis}
+
\operatorname{span}\left\{P_{\subentity}\left(\left.\phi_i\right|_{\subentity}\right)\,:\,l_i\in\left.\dualbasis\right|_{\closure{\subentity}}\right\}.
\label{eq:first-sum}\end{equation}
As (for all $l_i\in\left.\dualbasis\right|_{\closure{\subentity}}$) $P_{\subentity}\left(\left.\phi_i\right|_{\subentity}\right)\in\ctrace{\polyspace}{\subentity}{\dualbasis}$, we see due to \cref{lemma:zeroth} that 
\[\unctrace{\polyspace}{\subentity}{\dualbasis}
\cap
\operatorname{span}\left\{P_{\subentity}\left(\left.\phi_i\right|_{\subentity}\right)\,:\,l_i\in\left.\dualbasis\right|_{\closure{\subentity}}\right\}
=\emptyset,\] and so the sum in \cref{eq:first-sum} is disjoint, ie
\begin{align*}
\unctrace{\polyspace}{\subentity}{\dualbasis}\oplus\ctrace{\polyspace}{\subentity}{\dualbasis}
&=
\unctrace{\polyspace}{\subentity}{\dualbasis}
\oplus
\operatorname{span}\left\{P_{\subentity}\left(\left.\phi_i\right|_{\subentity}\right)\,:\,l_i\in\left.\dualbasis\right|_{\closure{\subentity}}\right\}.
\end{align*}
From this, we conclude that
\begin{align*}
\ctrace{\polyspace}{\subentity}{\dualbasis}
&=
\operatorname{span}\left\{P_{\subentity}\left(\left.\phi_i\right|_{\subentity}\right)\,:\,l_i\in\left.\dualbasis\right|_{\closure{\subentity}}\right\}.
\end{align*}
\end{proof}

\end{document}
